% Répartition des tâches du sous-groupe plateforme (lot L1), point de séance du 05/11/2026.
% Compiler depuis ce dossier : tectonic -o .. 2026-11-05_repartition-taches-plateforme.tex
\documentclass[10pt,a4paper]{article}
\usepackage{style_cr}
\hypersetup{pdftitle={Répartition des tâches plateforme, 05/11/2026}, pdfauthor={Reda El Mekkaoui},
            pdfsubject={Projet 20 Cerema, lot L1 Plateforme}}

\begin{document}

\entete{PROJET 20 CEREMA \quad·\quad LOT L1 PLATEFORME \quad·\quad POINT DE SÉANCE DU 05/11/2026}
       {Répartition des tâches plateforme}
       {Une dernière tentative sur le GPS, puis moteurs, télémétrie et ArduPilot pour piloter le bateau.}
       {Sous-groupe : Ewen, Rana, Selma}

\begin{enbref}
\begin{itemize}
  \item On garde un \textbf{GPS}, comme le souhaite Antoine Douchin (Cerema) : une dernière tentative avec le module actuel, sinon on en change.
  \item Chacun avance sur sa partie : \textbf{Ewen} le GPS et la télémétrie, \textbf{Rana} le branchement des moteurs et des ESC, \textbf{Selma} ArduPilot et QGroundControl.
  \item Une fois les tâches finies, \textbf{point commun} pour que tous comprennent le contrôleur de vol et le pilotage des moteurs.
\end{itemize}
\end{enbref}

\section{Ordre des étapes}

\begin{center}
\begin{tikzpicture}[
    etape/.style={rectangle, rounded corners=6pt, draw=#1!40, line width=0.5pt, fill=#1!6,
                  minimum width=5.1cm, minimum height=1.55cm, align=center, inner sep=4pt},
    num/.style={circle, fill=#1, text=white, font=\scriptsize\bfseries, minimum size=0.44cm, inner sep=0pt},
    quand/.style={font=\scriptsize\bfseries, text=#1},
    fl/.style={-{Stealth[length=5pt,width=4.5pt,round]}, draw=cExt, line width=0.8pt},
  ]
  \node[etape=cDec]  (a) at (0,0)    {\small\textbf{Chacun sa partie}\\[-1pt]{\footnotesize\color{muted}GPS, télémétrie, moteurs, ArduPilot}};
  \node[etape=cPlat] (b) at (5.75,0) {\small\textbf{Point commun}\\[-1pt]{\footnotesize\color{muted}contrôleur de vol, pilotage des moteurs}};
  \node[etape=cExt]  (c) at (11.5,0) {\small\textbf{Mission préprogrammée}\\[-1pt]{\footnotesize\color{muted}objectif du lot L1, par waypoints}};
  \draw[fl] (a) -- (b); \draw[fl] (b) -- (c);
  \node[quand=cDec,  above=3pt] at (a.north) {\tikz[baseline=(n.base)]\node[num=cDec](n){1};\enspace MAINTENANT};
  \node[quand=cPlat, above=3pt] at (b.north) {\tikz[baseline=(n.base)]\node[num=cPlat](n){2};\enspace ENSUITE};
  \node[quand=cExt,  above=3pt] at (c.north) {\tikz[baseline=(n.base)]\node[num=cExt](n){3};\enspace BUT};
\end{tikzpicture}
\end{center}

\section{Qui fait quoi}

\begin{minipage}[t]{0.32\linewidth}
\begin{personne}{cEwen}{EWEN}{GPS et télémétrie}
\tache{GPS : dernière tentative}{Trouver un câble JST-GH 10 broches, modifier certaines broches, brancher le GPS sur le port GPS1. Peu de temps à y passer.}
\tache{Télémétrie}{La liaison marche, mais comme un simple câble. Analyser ce qui passe entre les modules et ce qu'on peut en tirer.}
\end{personne}
\end{minipage}\hfill
\begin{minipage}[t]{0.32\linewidth}
\begin{personne}{cRana}{RANA}{branchement}
\tache{Moteurs et ESC}{Les brancher.}
\tache{Documentation}{Comment relier tous les composants, en tenant compte maintenant du moteur.}
\end{personne}
\end{minipage}\hfill
\begin{minipage}[t]{0.32\linewidth}
\begin{personne}{cSelma}{SELMA}{ArduPilot et QGroundControl}
\tache{Paramètres}{Les étudier depuis QGroundControl.}
\tache{Mission}{Comment programmer une mission. À avancer dès la prochaine séance, une fois les branchements faits.}
\end{personne}
\end{minipage}

\section{Le GPS : une dernière tentative}

\begin{center}
\begin{tikzpicture}[
    boite/.style={rectangle, rounded corners=6pt, draw=#1!40, line width=0.5pt, fill=#1!6,
                  minimum height=1.0cm, align=center, inner xsep=10pt, font=\small},
    fl/.style={-{Stealth[length=5pt,width=4.5pt,round]}, draw=cExt, line width=0.8pt, rounded corners=4pt},
    lab/.style={font=\scriptsize\bfseries, text=#1, above, inner sep=2pt},
  ]
  \node[boite=cEwen] (t) at (0,0) {\textbf{Module actuel (Neo M8P)}\\[-1pt]{\footnotesize\color{muted}câble JST-GH 10 broches, port GPS1}};
  \node[boite=cSelma, anchor=west, minimum width=4.9cm] (ok) at (6.3,0.68) {\textbf{On le garde}};
  \node[boite=wait, anchor=west, minimum width=4.9cm] (ko) at (6.3,-0.68) {\textbf{On change de module}\\[-1pt]{\footnotesize\color{muted}plan B vu le 21/10 : module F9P}};
  \coordinate (j) at ($(t.east)+(0.45,0)$);
  \draw[draw=cExt, line width=0.8pt] (t.east) -- (j);
  \draw[fl] (j) |- (ok.west) node[lab=cSelma, pos=0.78] {ça marche};
  \draw[fl] (j) |- (ko.west) node[lab=wait, pos=0.78] {ça ne marche pas};
\end{tikzpicture}
\end{center}
{\footnotesize\color{muted} Un module GPS prêté par un autre groupe a fonctionné : on ne peut pas le garder, mais on a vu comment le GPS doit se comporter. Passer encore du temps sur le module actuel ne fait pas l'unanimité : point à voir avec les tuteurs.\par}

\section{Pour la suite}
\begin{itemize}
  \item Quand chacun aura fini sa tâche : point commun du sous-groupe. Tout le monde doit comprendre comment fonctionne le contrôleur de vol, comment piloter les moteurs et quels retours on reçoit sur l'ordinateur.
\end{itemize}

\redige{Rédigé par Reda El Mekkaoui d'après la transcription automatique du point de séance. Corrections : « Eoen » → Ewen, « Ranel » → Rana, « l'Artois pilote » → ArduPilot, « GST, GH » → JST-GH. Si quelque chose est faux, dites-le, on corrige.}

\end{document}
